\subsection{Scelta dei compiti e degli scenari}
Al fine di condurre una valutazione di usabilità efficace e rappresentativa, la selezione dei compiti e degli scenari è stata effettuata seguendo un approccio mirato e non casuale. Le prove sono state progettate per esplorare le funzionalità principali e i flussi di navigazione più critici della piattaforma \textit{BrainBox}. L'obiettivo principale è stato quello di simulare l'esperienza reale di uno studente universitario alle prese con la gestione, la condivisione e l'elaborazione dei propri appunti, ponendo particolare attenzione all'interazione con gli strumenti innovativi basati sull'intelligenza artificiale, quali le funzioni di trascrizione e generazione audio.

Per garantire una valutazione che fosse completa e il più possibile naturale, lo svolgimento delle prove è stato diviso in due fasi successive, pensate per lasciare all'utente gradi di libertà differenti.

Nella prima fase è stato richiesto di eseguire \textbf{sei compiti} di difficoltà variabile, focalizzati sulle operazioni più comuni e fondamentali della piattaforma. Oltre a permetterci di raccogliere metriche precise sulle singole funzionalità, questo primo passaggio è servito anche come vero e proprio addestramento: portando a termine azioni specifiche e guidate, il partecipante ha potuto prendere confidenza con l'interfaccia e capirne le dinamiche di base.

Nella seconda fase sono stati invece proposti \textbf{due scenari} d'uso. A differenza del singolo compito, lo scenario ha lo scopo di inserire il partecipante in un contesto realistico, dandogli un obiettivo generale e verosimile. L'intento era quello di favorire l'immedesimazione del tester nella situazione descritta, lasciandogli piena libertà nello scegliere le azioni da compiere e i percorsi di navigazione da seguire per arrivare all'obiettivo prefissato.

In fase di progettazione dell'esperimento, sia i compiti che gli scenari sono stati accuratamente calibrati tenendo in considerazione i seguenti aspetti:
\begin{itemize}
    \item \textbf{Difficoltà crescente:} le prove sono state ordinate in modo da far prendere gradualmente confidenza all'utente con il sistema, partendo da azioni basilari per arrivare a flussi più complessi, al fine di metterlo a proprio agio fin dai primi minuti.
    \item \textbf{Durata limitata:} ogni prova è stata concepita per essere risolta in un lasso di tempo contenuto, stabilendo un limite massimo (10 minuti) per evitare di generare frustrazione, noia o un carico cognitivo eccessivo nel partecipante.
    \item \textbf{Rappresentatività:} i compiti selezionati rispecchiano le reali e più frequenti necessità di uno studente universitario sulla piattaforma.
\end{itemize}

È importante sottolineare che, per rispettare il vincolo della durata limitata dell'intero esperimento, non è stato possibile coprire l'assoluta totalità delle funzionalità offerte da \textit{BrainBox}, concentrando l'attenzione esclusivamente sui flussi più critici.

\subsubsection{Compiti}
Di seguito vengono descritti i sei compiti sottoposti agli utenti, formulati per guidarli passo dopo passo nell'esplorazione del sistema:

\begin{itemize}
    \item \textbf{Compito 1: Registrazione alla piattaforma.} È stato richiesto di creare un nuovo account inserendo una serie di dati specifici forniti in traccia (nome fittizio, email, password, Università degli Studi di Brescia, Dipartimento e Corso di studi). Lo scopo era testare l'accessibilità e la fluidità del form di registrazione iniziale.
    
    \item \textbf{Compito 2: Caricamento di un file in una cartella.} L'utente doveva creare una nuova cartella nominata "Archivio" all'interno della sezione \textit{Appunti} e caricarvi un file PDF preimpostato sul desktop (\textit{Appunti\_Lezione.pdf}).
    
    \item \textbf{Compito 3: Pianificare un esame.} Questo task mirava a testare la sezione di pianificazione di un esame: si richiedeva di creare l'esame "Architettura" impostando una data specifica, aggiungervi il task "Riassumere il primo capitolo" e allegarvi il documento PDF caricato nel compito precedente.
    
    \item \textbf{Compito 4: Rinomina un file.} Un'operazione di gestione file di base in cui si chiedeva di rinominare il file PDF precedente in "Riassunto Primo Semestre".
    
    \item \textbf{Compito 5: Salvataggio di un file della community.} Questo compito introduceva l'aspetto social della piattaforma: l'utente doveva cercare nella sezione \textit{Community} i documenti relativi al corso di "CyberSecurity", salvarne uno nella propria area personale e verificare la corretta riuscita dell'operazione.
    
    \item \textbf{Compito 6: Commentare un file della community.} In quest'ultimo compito si chiedeva di aprire un file pubblico della community e testare il modulo dei commenti inserendo la frase standard: "Documento molto utile, grazie!".
\end{itemize}

\subsubsection{Scenari}
L'esperimento è proseguito con la somministrazione di due scenari, caratterizzati da una narrazione più ampia e dall'assenza di istruzioni meccaniche \textit{step-by-step}.

\begin{itemize}
    \item \textbf{Scenario 1: "Il contributore attivo"}
        \begin{itemize}
            \item \underline{Contesto}: \textit{"Hai appena superato l'esame di Analisi 1 con un ottimo voto. Visto che in passato la community ti è stata molto d'aiuto, hai deciso di ricambiare il favore mettendo a disposizione di tutti i tuoi schemi."}
            \item \underline{Obiettivi richiesti}: L'utente doveva prima caricare un file PDF dal desktop sul proprio profilo, impostando correttamente la privacy su "Pubblico". Successivamente, doveva simulare la preparazione per un nuovo esame ("Fondamenti di Informatica"): gli è stato richiesto di creare l'esame nell'apposita sezione, cercare materiale utile nella community e collegarlo direttamente all'esame appena creato per averlo pronto allo studio.
        \end{itemize}
    \item \textbf{Scenario 2: "L'organizzatore seriale"}
        \begin{itemize}
            \item \underline{Contesto}: \textit{"È iniziato il nuovo anno accademico e la tua priorità è avere tutto il materiale ordinato in modo maniacale fin dal primo giorno, mantenendo il tuo spazio di lavoro rigorosamente privato."}
            \item \underline{Obiettivi richiesti}: In questo scenario l'utente doveva creare una nuova cartella dedicata all'esame di "Basi di Dati". Al suo interno doveva caricare un documento di testo (il \textit{Syllabus}). Infine, per testare le funzionalità avanzate di intelligenza artificiale, gli veniva richiesto di caricare all'interno della stessa cartella un file audio MP3 (simulando la registrazione di una lezione) e di utilizzare gli strumenti AI della piattaforma per convertirlo automaticamente in un file di testo testuale.
        \end{itemize}
\end{itemize}